\section{Two cones with two surviving Hom groups}
\label{cone:section}

We now use the two degree-three classes to construct an $A$-bimodule
$\mathcal C$, projective on each side, for which the stable maps from $X$
to $\mathcal C\otimes_A X[a]$ survive only in degrees $0$ and $3$.
The degree-three group will detect the maps constructed in the next section.
Throughout this section, put
\[
 A=T\otimes_{\kk}T,\qquad X=s\otimes_{\kk}s,\qquad R=A^e.
\]
The algebras $A$ and $R$ are symmetric: their symmetrizing forms are the
tensor products of the form on $T$ and, for opposite factors, the same form.
Write
\[
 H^a=\stHom{A}{X}{X[a]}.
\]
The preceding calculation gives
\begin{equation}
 H^{\geq0}=\kk[\tau_1,\tau_2],\qquad |\tau_1|=|\tau_2|=3.
 \label{cone:positive-ring}
\end{equation}
In degree zero this is also the stable endomorphism space: $X$ is one-dimensional,
and its nonzero degree-three self-extension proves that it is nonprojective.

\subsection{Negative degrees and multiplication}

We first record the symmetric stable duality, including the compatibility
with multiplication that will be needed below; see also
\cite[Section~2]{Linckelmann2013}.
For finite modules $M,N$ over a symmetric algebra $G$, it says that
\begin{equation}
 D\stHom{G}{M}{N}\simeq\stHom{G}{N}{M[-1]},
 \qquad D=\Hom_{\kk}(-,\kk),
 \label{cone:duality}
\end{equation}
naturally in $N$.
Here is a direct proof with our indexing conventions. Choose a complete
resolution $Q$ of $M$, so that $M=\coker(Q_1\to Q_0)$.
For a finite projective $Q_i$, symmetry gives a natural isomorphism
\[
 D\Hom_G(Q_i,N)\simeq\Hom_G(N,Q_i).
\]
It suffices to check this for $Q_i=G$, where it is the isomorphism
$DG\simeq G$, and then to pass to finite sums and summands.
Dualizing the cochain complex $\Hom_G(Q,N)$ therefore gives the chain
complex $\Hom_G(N,Q)$. Its homology at zero consists of maps into
$\ker(Q_0\to Q_{-1})=\Omega M$, modulo maps through $Q_1$.
These are exactly the projective factorizations, since maps from a projective
to $\Omega M$ lift through $Q_1\twoheadrightarrow\Omega M$.
Vector-space duality is exact, proving \eqref{cone:duality} and its naturality.

In particular,
\begin{equation}
 DH^a\simeq H^{-a-1}.
\label{cone:negative-duality}
\end{equation}
For $m\geq0$, let $\mathcal P_m$ denote the degree-$m$ part of the polynomial ring
$\kk[\tau_1,\tau_2]$ when each variable is assigned degree one.
Thus
\begin{equation}
 H^{3m}=\mathcal P_m,\qquad H^{-3m-1}=D\mathcal P_m\quad(m\geq0),
 \label{cone:support}
\end{equation}
and all other groups vanish.

These identifications also determine the action of the two classes in
negative degrees. Postcomposition by $\tau_i[a]$ gives
$H^a\to H^{a+3}$. By naturality of \eqref{cone:duality}, its transpose,
after shifting the source back to $X$, is
\[
 H^{-a-4}\longrightarrow H^{-a-1},\qquad
 g\longmapsto g[3]\circ\tau_i.
\]
When $a\leq-4$, the last map is multiplication by $\tau_i$ in the
polynomial ring \eqref{cone:positive-ring}. The action within negative
degrees is therefore the dual of polynomial multiplication.
There is no nonzero product crossing from negative to nonnegative degree:
the only nonzero possible source is $H^{-1}$, and its target is $H^2=0$.

\subsection{Finite bimodules from projective tails}

We will form the two cones as ordinary complexes before passing to stable
bimodules. We give the passage explicitly to retain finite representatives
and exact tensor evaluation.

Bimodules attached to Hochschild classes and their iterated tensor products
have a precursor in the support-variety constructions of Erdmann,
Holloway, Snashall, Solberg and Taillefer \cite[Section~2]{ErdmannEtAl2003}.
Here we compute the full stable Hom profile directly, including negative
degrees and the top attaching map.

If $V$ is a bounded below complex of finite projective $R$-modules that is
exact in sufficiently high degrees, put
\begin{equation}
 \langle V\rangle=
 \bigl(\coker(V_{N+1}\longrightarrow V_N)\bigr)[N]
 \label{cone:tail-definition}
\end{equation}
in the stable category, with $N$ sufficiently large.
This does not depend on $N$: writing $V_N^{\mathrm c}$ for this cokernel,
exactness gives
\[
 0\longrightarrow V_{N+1}^{\mathrm c}
 \longrightarrow V_N\longrightarrow V_N^{\mathrm c}
 \longrightarrow0.
\]
Thus $V_{N+1}^{\mathrm c}\simeq V_N^{\mathrm c}[-1]$ stably.
Chain maps defined in sufficiently high degrees induce maps on these
cokernels, compatibly with increasing $N$ and with shifts. A chain homotopy
changes the bottom induced map by a factorization through a projective,
so homotopic tail maps induce the same stable map.

For a degreewise split short exact sequence of such complexes
$0\to B\to V\to D\to0$, their cokernels at $N$ form a short exact
sequence once $H_{N+1}(D)=0$.
Indeed, a possible extra relation in the $B$-cokernel is obtained by lifting
a cycle in $D_{N+1}$. That cycle is a boundary; subtracting a lifted
boundary leaves a relation already in $B$.
The resulting stable sequence has the usual long exact stable Hom sequence:
apply $\Hom$ from a complete projective resolution to the short exact
sequence and use its cohomology description of stable Hom.
If the cross term of the differential is
$r:D_j\to B_{j-1}$, the connecting stable map
$\langle D\rangle\to\langle B\rangle[1]$ is induced by $r$.
This follows by lifting along the degreewise splitting and taking the
$B$-component of the differential, which is precisely $r$.

Call a bimodule \emph{side-projective} if it is projective separately as a
left and as a right $A$-module. Finite stable shifts of a finite
side-projective bimodule have finite side-projective representatives.
For a syzygy, take a finite $R$-projective surjection; its kernel sequence
splits on each side because its quotient is side-projective.
For a cosyzygy, take a finite $R$-projective-injective embedding; its
cokernel sequence splits on each side because its submodule is
side-injective, by symmetry of $A$.
Likewise a stable map $M\to N$ between such bimodules has a finite
side-projective cone: represent the map by $f$, embed $M$ into a finite
$R$-projective $I$, and take the cokernel of
$(f,i):M\to N\oplus I$. Again the sequence splits on each side.

\begin{lemma}
\label{cone:finite}
Let $P$ be the bimodule bar resolution of $T$ and let
$p:P\to P[3]$ be the selected comparison map. Form the ordinary complex
\[
 K=\operatorname{Cone}(p:P[-3]\to P)
       \otimes_{\kk}\operatorname{Cone}(p:P[-3]\to P).
\]
Then $K$ has finite projective $R$-terms and is exact in degrees greater
than zero. For every $N\geq1$, its cokernel
$C_N=\coker(K_{N+1}\to K_N)$ is finite and side-projective.
Consequently $\mathcal C=\langle K\rangle$ has a finite side-projective
module representative. The same conclusions hold for the first two
filtration levels of $K$.
\end{lemma}

\begin{proof}
Write $P^{\mathrm{tot}}=P\otimes_{\kk}P$.
In each total degree there are only finitely many summands, each an
external tensor of projective bimodules; these are projective over $R$.
The complex $P^{\mathrm{tot}}$ resolves $A$.
Label the four cone cells by subsets $S\subseteq\{1,2\}$, where $i\in S$
means that the shifted cell is selected in cone factor $i$.
The $S$-cell is $P^{\mathrm{tot}}[-2|S|]$; its differential to the cell
$S\setminus\{i\}$ is the shifted comparison map $p_i$.
Filtering by $|S|$ gives the three associated levels
\begin{equation}
 P^{\mathrm{tot}},\qquad
 (P^{\mathrm{tot}}[-2])^{\oplus2},\qquad
 P^{\mathrm{tot}}[-4].
 \label{cone:cells}
\end{equation}
The cell homology is supported at $0,-2,-4$, respectively. The filtration
long exact sequences therefore show that $K$ is exact in positive degrees.

For $N\geq1$, the cokernel filtration is exact by the preceding observation.
Its factors, with multiplicities $1,2,1$, are
\[
 \coker\bigl(P^{\mathrm{tot}}_{N+2r+1}
       \to P^{\mathrm{tot}}_{N+2r}\bigr),\qquad r=0,1,2.
\]
The augmented resolution of $A$ splits on either $A$-side, since the
regular module is projective; inductively all its syzygies are projective
on each side. Thus these filtration factors are side-projective.
Their extensions split on each side, proving the same property for $C_N$.
Finally $C_N[N]$ has a finite side-projective representative by the shift
construction above. In particular one may take $N=1$ and one finite
cosyzygy of $C_1$.
\end{proof}

Forming the external tensor of ordinary complexes gives all four cells
of \eqref{cone:cells} as projective resolutions of shifts of the regular
$A$-bimodule. We pass to their stable representatives only after this
construction.

Evaluation on a finite left module means tensoring over $A$ on the right.
Tensoring with a side-projective bimodule is an exact functor of left
$A$-modules and preserves projectives.
Every $R$-projective evaluates to an $A$-projective, since a free bimodule
evaluates to a module of the form $A\otimes_{\kk}V$.
For evaluation at a fixed module, the right-split sequences used above
remain exact. Their evaluations therefore preserve the stable shifts and
long exact Hom sequences in the construction.
Accordingly the filtration \eqref{cone:cells}, after evaluation on $X$,
has stable cells
\[
 X,\qquad X[-2]^{\oplus2},\qquad X[-4].
\]
Its adjacent connecting maps are the shifted actions of $\tau_1,\tau_2$:
they are induced by the original bimodule comparison maps tensored with
$X$. This specifies the attaching maps whose Hom groups we now compute.

\subsection{The Hom profile and its top projection}

\begin{proposition}
\label{cone:profile}
For the finite side-projective bimodule $\mathcal C$ of
Lemma~\ref{cone:finite}, put
\[
 W^a=\stHom{A}{X}{(\mathcal C\otimes_A X)[a]}.
\]
For every integer $a$,
\[
 W^a=\begin{cases}\kk,&a=0\text{ or }3,\\0,&a\notin\{0,3\}.
 \end{cases}
\]
Moreover the projection to the top cell induces an isomorphism
\begin{equation}
 W^3\xrightarrow{\ \sim\ }H^{-1}.
 \label{cone:top-projection}
\end{equation}
\end{proposition}

\begin{proof}
Let $\mathcal C_{\leq1}$ denote the first two filtration levels, represented
by their projective tails as above. In total cohomological degree $a$,
the Hom groups of the three associated levels are
\[
 B^a=H^a,\qquad D^a=(H^{a-2})^2,\qquad E^a=H^{a-4}.
\]
The adjacent connecting maps are
\[
 d_a:D^a\longrightarrow B^{a+1},\qquad
 d_a(u,v)=\tau_1u+\tau_2v,
\]
\[
 d'_a:E^a\longrightarrow D^{a+1},\qquad
 d'_a(z)=(\tau_2z,\tau_1z).
\]
The order in the column records that deleting the first slot from the
full cell leaves the second singleton, and conversely.

We need the degreewise exact Koszul sequence
\begin{equation}
 0\longrightarrow \mathcal P_{m-1}
 \xrightarrow{(\tau_2,\tau_1)^t}\mathcal P_m^{\oplus2}
 \xrightarrow{(\tau_1,\tau_2)}\mathcal P_{m+1}
 \longrightarrow0\qquad(m\geq0),
 \label{cone:koszul}
\end{equation}
where $\mathcal P_j=0$ for $j<0$.
The row is surjective because every positive-degree monomial is divisible
by one of the variables. The column is injective, and coprimality of the
variables shows that every relation $\tau_1u+\tau_2v=0$ is of the form
$(u,v)=(\tau_2w,\tau_1w)$.
On negative groups, dualize \eqref{cone:koszul} and interchange its two
middle components. By the multiplication calculation above this gives
exactly the sequence through $D^a$ there. At the seam its terminal zero
space is the actual group $H^2=0$. We obtain, in all degrees,
\begin{align}
 \ker d_a&=\im d'_{a-1},\nonumber\\
 \coker d_{a-1}&=
 \begin{cases}\kk,&a=0,\\0,&a\ne0,\end{cases}
 &\ker d'_a&=
 \begin{cases}\kk=H^{-1},&a=3,\\0,&a\ne3.
 \end{cases}
 \label{cone:koszul-consequences}
\end{align}
For the column kernel, its maps on $D\mathcal P_m$ for $m\geq1$ are dual to the
surjective positive row and hence injective. Its remaining source is
$D\mathcal P_0=H^{-1}$, whose two products vanish in $H^2$.

It remains to check that the actual filtration has no further effect on
these scalar classes. Put
\[
 K_1^a=\stHom{A}{X}{(\mathcal C_{\leq1}\otimes_A X)[a]}.
\]
The first filtration long exact sequence gives
\begin{equation}
 0\longrightarrow\coker d_{a-1}\longrightarrow K_1^a
 \xrightarrow{q_a}\ker d_a\longrightarrow0.
 \label{cone:first-les}
\end{equation}
The top-cell connecting map is $r_a:E^a\to K_1^{a+1}$, and its projection
satisfies $q_{a+1}r_a=d'_a$. Thus $\ker r_a\subseteq\ker d'_a$, so it can
be nonzero only when $a=3$. At that index $K_1^4=0$ by
\eqref{cone:first-les}, since $D^4=(H^2)^2=0$. Hence
$\ker r_3=E^3=H^{-1}$.

The image of $r_{a-1}$ surjects onto $\ker d_a$ by
\eqref{cone:koszul-consequences}. When $a\ne0$, the left term of
\eqref{cone:first-les} vanishes, so $r_{a-1}$ is surjective.
When $a=0$, that sequence gives $K_1^0=\kk$, whereas the only incoming
source is $E^{-1}=H^{-5}=0$. Therefore $\coker r_{a-1}$ is $\kk$ exactly
at $a=0$ and zero elsewhere.

Finally the second filtration long exact sequence gives
\[
 0\longrightarrow\coker r_{a-1}\longrightarrow W^a
 \longrightarrow\ker r_a\longrightarrow0.
\]
This proves the asserted groups. At $a=3$ its left term vanishes and its
right term is all of $E^3=H^{-1}$, so the actual top-cell projection is
an isomorphism, proving \eqref{cone:top-projection}.
\end{proof}

The two exact sequences in the proof account for the complete filtration,
including any attaching map not determined by its adjacent component.
We have therefore obtained the required two-scalar profile together with
a concrete way to detect its degree-three class: after shifting the cone
by three, it suffices to project to $X[-1]$. The next construction supplies
maps with a nonzero projection there.
